\section{Discussion}
\label{sec:discussion}

% =============================================================================
\subsection{What should be reported}
\label{sec:discussion:reporting}

The main lesson is methodological. All-site fault metrics are meaningful
only together with the fault-site population they average over. Reporting
\AvgDP{} without the representation granularity can make a decomposition
look like a defense or a regression depending on the chosen baseline. A
minimum report should therefore include the number of fault sites,
fault sites per primary output, whether sites are all internal gates or a
targeted layer, and the synthesis pass that created the baseline.

For cross-representation comparisons, we recommend pairing all-site
metrics with at least one controlled metric. A decomposed-AIG baseline at comparable small-gate granularity
tests whether the effect is specific to the logic primitive rather than
to decomposition size. A target-layer metric tests whether the
attacker-facing output layer changes when interior nodes are excluded.
In our benchmark, both controls overturn the naive LUT-to-MIG
interpretation.

% =============================================================================
\subsection{What remains useful}
\label{sec:discussion:robust}

The granularity control does not make the artifact negative. It separates
three kinds of claims. First, the LUT-to-MIG \AvgDP{} reduction is a
real property of that pair of netlists, useful if a designer specifically
cares about all internal sites in those representations. Second, the
target-layer result shows that this reduction should not be promoted to a
representation-independent fault-resistance claim. Third, the ARX
\AvgBW{} result survives the controls and is therefore a stronger
design fact.

% =============================================================================
\subsection{Limits}
\label{sec:discussion:limits}

\textit{Fault model.} We model one transient \bitflip{} at one
gate-output wire. Stuck-at faults, delay faults, multi-bit faults, and
physical injection footprints may interact with representation
granularity differently.

\textit{Input model.} For large circuits we use a fixed 10,000-vector
uniform random sample; small circuits are exhaustive. An adversary who
can choose inputs may target vectors that maximize or minimize
propagation. Worst-case input search is a separate attack problem.

\textit{Target layer.} Output-driving gates are a clean controlled layer,
not a complete physical adversary model. They are useful because the
layer exists in every representation and avoids the interior-node
population shift, but other fixed layers could be studied.

\textit{Sample size.} The ARX result uses three ARX primitives and seven
substitution-based primitives. Its no-overlap pattern is strong in this
benchmark, but a wider set of addition-chain and substitution designs is
needed before treating the ratio as a general law.

\textit{End-to-end DFA.} None of the metrics directly measures key
recovery. \FCpct{}, \AvgDP{}, and \AvgBW{} describe observability and
output corruption. Algebraic usefulness of those differences remains
cipher-specific.
